\begin{textsample}{POS Dim 1 – global\_north\_2024\_04 – Score 38.00 – global\_north\_2024\_04/107142833.txt}  \label{ex:f1_pos_198}
Dominic Allen, the head of the United Nations \textbf{Population} Fund ( UNFPA ) covering Gaza, said he was"terrified"of what could happen if the war went on any longer.
He told AFP the situation was"beyond catastrophic"with gaunt and starving people spending their days searching for \textbf{food}, and \textbf{medicine} running desperately low.
The British-born official, who spent a week in the \textbf{besieged} Palestinian \textbf{territory} last month, said even when \textbf{aid} got through the border, there were still major problems getting it to those who need it most, particularly women and girls.
QUESTION : What is Gaza like now?
ANSWER :"What I saw across the Gaza Strip is beyond catastrophic.
I have been to Gaza many times before this war, and what I saw ( this time ) was truly heartbreaking.
Gaza is dust."Everybody we drove by and many we spoke to were gaunt, starving, \textbf{hungry} and looking frail.
Everybody was looking for \textbf{food}."We ' re really worried about pregnant women and breastfeeding @ @ @ @ @ @ @ @ @ @ ( the only one functioning in the north ) say women are giving birth to smaller babies caused by \textbf{malnutrition}, \textbf{dehydration} and fear."They ' re telling us anecdotally that they are seeing no normal-sized babies being delivered in Gaza, an increase in the number of stillborns and neo-natal deaths as well."\textbf{Delivery} rooms are overwhelmed.
A midwife described women giving birth on the floor because they were at maximum capacity."They ' re having to use thread for umbilical ties."QUESTION : Many international organisations have raised concerns about how difficult it is to get \textbf{aid} into Gaza.
How has it been for you?
ANSWER :"UNFPA has had a number of our suppliers denied entry into Gaza at checkpoints and screenings."Access to the north is very challenging... many UN missions to the north have been denied over past months."People are on the verge of \textbf{famine} right now in Gaza.
This is caused by a huge backlog in \textbf{supplies} and assistance. @ @ @ @ @ @ @ @ @ @ ANSWER :"We brought what they needed the most, the anaesthetics, the oxytocin ( and other items that need to be kept \textbf{cold} ) and put them in the back of our armoured vehicle and deliver them by hand to the \textbf{hospital}."QUESTION : What do the \textbf{hospitals} most need?
ANSWER :"They ' re asking for \textbf{fuel}.
Their \textbf{hospitals} are running on fumes.
One said if a \textbf{patient} needs \textbf{surgery}, they have to take a canister of gasoline or diesel ( with them ) to run the generator in the operating theatre.
They lack some of the most basic items to help support and give safe birth care for mothers and babies."QUESTION : Are you able to deliver period products for women and girls?
ANSWER :"There are many thousands of cases of insufficient menstrual hygiene items, really shocking reports of women who have had to fashion their own sanitary products from pieces of \textbf{tents}.
So one of the UNFPA ' s priorities is to get in at @ @ @ @ @ @ @ @ @ @ kits."QUESTION : Are you afraid of what will happen, particularly if there is an \textbf{offensive} on Rafah?
ANSWER :"Every person I spoke to in Rafah is so fearful about what will come next for a potential \textbf{ground} incursion in Rafah ( in the south )."I left with a real sense of fear about what could come next.
If the Emirati \textbf{Hospital} is managing nearly 100 births every single day, what will happen to those pregnant women and babies?"What will happen to the 1.2 million people currently living and \textbf{sheltering} in Rafah right now?
A \textbf{city} which only housed 250,000 before is bursting in capacity."I left Gaza terrified of what could come next for them.
When you see small, frail babies cuddled together because there are not enough \textbf{incubators}, two or three in one, the frailty of life is so clear."The unprecedented Hamas attack on October 7 resulted in about 1,160 deaths in Israel, mostly \textbf{civilians}, according @ @ @ @ @ @ @ @ @ @ \textbf{Disclaimer} : The content of this article is \textbf{syndicated} or provided to this website from an \textbf{external} \textbf{third} party \textbf{provider}.
We are not responsible for, and do not control, such \textbf{external} websites, \textbf{entities}, applications or media \textbf{publishers}.
The body of the \textbf{text} is provided on an"as is"and"as \textbf{available}"\textbf{basis} and has not been \textbf{edited} in any way.
Neither we nor our \textbf{affiliates} \textbf{guarantee} the \textbf{accuracy} of or endorse the views or opinions expressed in this article.
Read our full \textbf{disclaimer} policy here.

% matched lemmas: accuracy, affiliate, aid, available, basis, besieged, city, civilian, cold, dehydration, delivery, disclaimer, edit, entity, external, famine, food, fuel, ground, guarantee, hospital, hungry, incubator, malnutrition, medicine, offensive, patient, population, provider, publisher, shelter, supply, surgery, syndicate, tent, territory, text, third
\end{textsample}
